% !TeX root = ../MAgent_Theory_2.0.tex
\chapter{State-Controlled Erasure and Persistent Groups}
\label{ch:smooth-erasure}

There is a conditional first-stage construction inside the original smooth KL potential. A prescribed channel can lose its dependence on some internal laws when a dynamically evolving gate coordinate enters a flat part of that channel. The remaining agents can be distinct, their preferred groups can change, and the retained dynamics can then become exactly autonomous. The mechanism is an engineered constitutive choice. It supplies neither a general explanation of sufficient interfaces nor a second formation stage. The general laws of the preceding chapter remain the framework; the existence example below uses an explicitly invariant family of positive finite laws.

\section{A homogeneous population with internal gate coordinates}

Write $b_x(k)=(1+kx)/2$ on $\{-1,1\}$, set $d(x,z)=\KL(b_x\Vert b_z)$, and put $\psi(x)=d(x,0)$. Four agents carry
\begin{equation}
 q_i=b_{x_i}\otimes b_{z_i}\otimes b_{h_i},\qquad
 s_i=b_{y_i}\otimes b_{w_i}\otimes b_{\widehat h_i}.
 \label{erase:family}
\end{equation}
The first coordinate is public, the second private, and the third controls a channel. Their unit Fisher metric has diagonal entries $1/(1-u^2)$ for each mean $u$. The product preparation within each sector is an assumption. It is invariant under the full table dynamics specified below, rather than a projection imposed after each step.

Take $J=\log3$ and $h_*=-1/2$. Use the proper reference $r=b_0\otimes b_0\otimes b_{h_*}$ and normalized evaluator
\begin{equation}
 G_J(k\mid m)=\frac{e^{Jkm}}{2\cosh J}
 \label{erase:evaluator}
\end{equation}
on the public coordinates, with independent private output $b_0$ and gate output $b_{h_*}$. The local term is exactly
\begin{equation}
 \begin{split}
 E_{0,i}={}&\psi(x_i)+\psi(y_i)-Jx_iy_i+\log\cosh J\\
 &+\psi(z_i)+\psi(w_i)+d(h_i,h_*)+d(\widehat h_i,h_*).
 \end{split}
 \label{erase:local}
\end{equation}
All agents have the same reference and evaluator. Every nonself agent is an allowed source, with uniform prior and fixed temperature $T>0$. We use the exact Gibbs-reduced version of the original potential. The bits are reference coordinates: a change of presentation carries the reference laws, channels, and readout maps with it, so the KL and Fisher covariance statements apply.

Define $\chi(r)=0$ for $r\leq0$, $\chi(r)=1$ for $r\geq1$, and, for $0<r<1$,
\begin{equation}
 \chi(r)=\frac{e^{-1/r}}{e^{-1/r}+e^{-1/(1-r)}}.
 \label{erase:flat-step}
\end{equation}
This increasing function is smooth, all its derivatives match at the endpoints, and $M_\chi=\max|\chi'|<\infty$. Fix $h_0=1/4$ and $\delta>0$, and set
\begin{equation}
 H_q=\frac14\sum_i h_i,\quad
 H_s=\frac14\sum_i\widehat h_i,\quad
 \alpha(H)=\chi((H-h_0)/\delta).
 \label{erase:gate}
\end{equation}
The gate is a nonlocal function of internal laws. Its threshold and flat region are constitutive data. There is no external time schedule.

Choose gains $0<\varepsilon,\varrho<1$. On a belief edge $j\to i$, let the public output kernel be $(1-\varepsilon)b_{x_i}+\varepsilon\delta_{k_{\rm pub}}$, and independently draw private output from $b_{\alpha(H_q)\varrho z_j}$ and gate output from $b_{h_*}$. Its compared law is
\begin{equation}
 \nu_{ij}^q=b_{(1-\varepsilon)x_i+\varepsilon x_j}
              \otimes b_{\alpha(H_q)\varrho z_j}\otimes b_{h_*}.
 \label{erase:compared-law}
\end{equation}
This normalized channel is strictly positive on the interior and depends smoothly on the current marginal means. It transmits public samples and uses the source's private law. The model channel exchanges $(x,z,h)$ with $(y,w,\widehat h)$. Define
\begin{equation}
 D_\varepsilon(x,v)=d(x,(1-\varepsilon)x+\varepsilon v).
\end{equation}
The original peer cost becomes
\begin{equation}
 D_{ij}^q=D_\varepsilon(x_i,x_j)
       +d(z_i,\alpha(H_q)\varrho z_j)+d(h_i,h_*).
 \label{erase:peer-cost}
\end{equation}
No fusion cost has been added.

The definitions extend to arbitrary positive tables on each three-bit space by taking the indicated marginal means. At a product table, the derivative of the receiver entropy is additive in its coordinates. Each compared law is a product, and every channel derivative is a function of a single coordinate because the channel parameters depend on marginal means. The local evaluator also gives additive scores within each sector. The full Fisher score is consequently a sum of coordinate scores on \cref{erase:family}. Differentiating a product shows that the full table flow preserves this family and induces the stated Bernoulli Fisher metric. This argument includes incoming and channel derivatives.

\section{Finite entry into an exactly erasing region}

\lemmaheading{Permanent gate closure under the original flow}{erase:gate-entry}
For sufficiently small $\delta>0$ and then sufficiently small $\varrho>0$ relative to $\delta$, there is an open class within \cref{erase:family}, with
\begin{equation}
 |z_i|,|w_i|<p=3/4,\qquad
 h_i,\widehat h_i\in(h_0+\delta/2,h_0+3\delta/4),
 \label{erase:initial-gates}
\end{equation}
for which both gates reach zero in finite time and remain zero.
\begin{proof}
The exact envelope derivative gives
\begin{equation}
 \dot h_i=-(1-h_i^2)\left[
 2(\operatorname{atanh}h_i-\operatorname{atanh}h_*)+
 \frac{\alpha'(H_q)}4\sum_{a,j}\beta_{aj}^*
       \partial_\alpha d(z_a,\alpha\varrho z_j)\right].
 \label{erase:gate-flow}
\end{equation}
The factor two comes from the local term and the receiver's source-independent gate cost. For $|z_a|,|z_j|\leq p$ and $\varrho\leq1/2$,
\begin{equation}
 \partial_\alpha d(z_a,\alpha\varrho z_j)
 =\varrho z_j\frac{\alpha\varrho z_j-z_a}
                       {1-\alpha^2\varrho^2z_j^2}.
 \label{erase:gate-derivative}
\end{equation}
Its absolute value is bounded by $B\varrho$, where one may take
$B=3p^2/[2(1-p^2/4)]$. The four row sums equal one, so the correction to the gate force is bounded by $B\varrho M_\chi/\delta$.

On $K=[1/8,1/2]$, define the positive constants
\begin{equation}
 \begin{aligned}
 m&=\min_{h\in K}2(1-h^2)
                (\operatorname{atanh}h-\operatorname{atanh}h_*),\\
 M&=\max_{h\in K}2(1-h^2)
                (\operatorname{atanh}h-\operatorname{atanh}h_*),
 \qquad V=M+m/2.
 \end{aligned}
\end{equation}
Choose $B\varrho M_\chi/\delta\leq m/2$. While its coordinates stay in $K$, each gate coordinate has $\dot h_i\leq-m/2$ and $|\dot h_i|\leq V$. Also choose $\delta$ so that the initial interval lies in $K$ and
\begin{equation}
 \frac{3V\delta}{2m}<h_0-1/8.
 \label{erase:time-choice}
\end{equation}
The sector mean reaches $h_0$ by time $3\delta/(2m)$, before any coordinate can leave $K$. On $H_q=h_0$, both $\alpha$ and $\alpha'$ vanish. Thereafter
\begin{equation}
 \dot h_i=-2(1-h_i^2)
           (\operatorname{atanh}h_i-\operatorname{atanh}h_*).
 \label{erase:closed-clock}
\end{equation}
All these coordinates start above $h_*$ and decrease toward it. Their mean therefore cannot leave the closed-gate region. The model-sector proof is identical.

It remains to justify the private bounds used in the argument. At $\varrho=0$ the private field is
$\dot z=-2(1-z^2)\operatorname{atanh}z$, strictly inward at $z=\pm p$. Private derivatives of \cref{erase:peer-cost} converge uniformly to this field as $\varrho\to0$, uniformly over normalized source rows and $0\leq\alpha\leq1$. The incoming terms are of order $\varrho$; differentiating a private mean does not differentiate $\alpha$. Compactness therefore preserves the strict inward signs for sufficiently small $\varrho$. The same holds for $w$. This closes the interval argument.
\end{proof}

Let $\tau$ be the later of the two gate-closure times. The construction uses a smooth but nonanalytic channel: its exactly flat region makes finite entry into exact erasure possible. Decreasing a small nonzero analytic gate would not establish this conclusion.

\section{A change of preferred groups without synchronization}

\theoremheading{Persistent heterogeneous pairs}{erase:pairs}
For sufficiently small $\varepsilon>0$, admit two public states in $(l,u)^2$ and two in $(-u,-l)^2$, with $l=0.49$, $u=0.51$, including every assignment of these signs to the four labels. These boxes are forward invariant. After $\tau$, each agent's unique preferred source in both sectors is the other agent in its box. There are initial private states for which the preferred pairs are instead crossed between these boxes before closure. The public states within each final pair can remain distinct at every finite time.
\begin{proof}
For the local public flow, the inequalities
\begin{equation}
 Jl-\operatorname{atanh}l>0,\qquad
 \operatorname{atanh}u-Ju>0
 \label{erase:well-margins}
\end{equation}
give strict inward signs on every face of the two boxes. They follow by differentiating $Jv-\operatorname{atanh}v$ near its zero $v=1/2$. Set $c=(1-\varepsilon)x+\varepsilon v$. Direct differentiation gives
\begin{equation}
 \begin{aligned}
 \partial_1D_\varepsilon(x,v)
 &=\operatorname{atanh}x-\operatorname{atanh}c
       +(1-\varepsilon)\frac{c-x}{1-c^2},\\
 \partial_2D_\varepsilon(x,v)
 &=\varepsilon\frac{c-x}{1-c^2}.
 \end{aligned}
 \label{erase:public-derivatives}
\end{equation}
Both are uniformly $O(\varepsilon^2)$ on $[-u,u]^2$; the first-order terms cancel in the first line. The full public force has one outgoing row and three incoming sources, each weighted by a number in $[0,1]$. A sufficiently small $\varepsilon$ preserves every strict inward margin, independently of the private state and the gate.

After gate closure, the private and gate costs in a row do not depend on its source, so Gibbs ranking uses $D_\varepsilon$ alone. Taylor's theorem in the first argument of $d$ gives
\begin{equation}
 \frac12(x-c)^2\leq d(x,c)
 \leq\frac{(x-c)^2}{2(1-u^2)}.
 \label{erase:cost-comparison}
\end{equation}
A same-box source has separation at most $u-l$, whereas an opposite-box source has separation at least $2l$. Since $(u-l)^2/(1-u^2)<(2l)^2$, the same-box partner is the unique preferred source throughout the invariant region. Cross-pair weights remain positive.

For an explicit initial reorganization, choose private means $\pm\zeta$, with $0<\zeta<p$, in both sectors, assigning each same-private-sign pair one positive and one negative public member. Initially $\alpha\geq1/2$. The private advantage of the same-sign partner over an opposite-sign source is
\begin{equation}
 d(\zeta,-\alpha\varrho\zeta)
 -d(\zeta,\alpha\varrho\zeta)
 =2\zeta\operatorname{atanh}(\alpha\varrho\zeta)
 \geq\varrho\zeta^2.
 \label{erase:crossed-pair-gap}
\end{equation}
The public cost range is at most $2u^2\varepsilon^2/(1-u^2)$. Choosing this smaller than $\varrho\zeta^2$ makes the initial preferred pairs the crossed pairs. Strict inequalities persist under small perturbations of the initial private values. Thus the final persistent grouping is not an unchanged initial source assignment.

Choose initial public states within a box separated by a fixed positive distance. Their speeds are uniformly bounded, whereas $\tau\leq3\delta/(2m)$. Taking $\delta$ sufficiently small keeps them distinct through $\tau$. Afterward the public field is autonomous, smooth, and permutation equivariant. Its equality diagonal for a pair is invariant. Backward uniqueness precludes first entry into that diagonal at a finite time. This proves finite-time distinctness; it does not give a uniform lower bound on their separation as $t\to\infty$.
\end{proof}

The parameter choices are compatible: first choose a sufficiently small $\delta$, then $\varrho$ sufficiently small relative to $\delta$, and finally $\varepsilon$ sufficiently small that $\varepsilon^2$ is small relative to $\varrho$ and the fixed public inward margins. All means stay in compact subsets of $(-1,1)$, giving global continuation in this admitted family.

\section{The exact inherited state}

For $t\geq\tau$ the reduced potential separates exactly:
\begin{equation}
 \begin{split}
 \mathcal E_{\rm red}={}&\mathcal E_{\rm pub}(\mathbf x,\mathbf y)\\
 &+2\sum_i[\psi(z_i)+\psi(w_i)
                +d(h_i,h_*)+d(\widehat h_i,h_*)].
 \end{split}
 \label{erase:separated-energy}
\end{equation}
Here $\mathcal E_{\rm pub}$ is the original coupled local KL and Gibbs-reduced peer potential with costs $D_\varepsilon$. The product Fisher metric separates as well. Discarding the private and gate laws therefore gives an exact autonomous public system without fitting another potential or mobility.

This independence is generally absent before gate closure. To see that this is an actual change of the vector field, expand one reduced row with $\alpha>0$ fixed. If
$A_{ij}=(x_j-x_i)^2/[2(1-x_i^2)]$, its first mixed public-private term is
\begin{equation}
 \frac{\alpha\varrho\varepsilon^2}{T}
 z_i\operatorname{Cov}_{j\ne i}(A_{ij},z_j).
 \label{erase:mixed-term}
\end{equation}
This follows by expanding $D_\varepsilon=\varepsilon^2A+O(\varepsilon^3)$, $d(z_i,\alpha\varrho z_j)=\psi(z_i)-\alpha\varrho z_i z_j+O(\varrho^2)$, and the logarithm of the row normalization. The covariance is under the uniform source prior. It is not identically zero even within the crossed-preference preparation. Fix $v=(-b,-a,a,b)$ with $l<a<b<u$, keep $z=\zeta(1,-1,1,-1)$ fixed, and vary only public $x_i=r v_i$ near $r=1$. The summed coefficient is
\begin{equation}
 \frac{2\zeta^2r^2}{9}\left[
 \frac{a^2-3b^2+6ab}{1-r^2b^2}
 +\frac{b^2-3a^2+6ab}{1-r^2a^2}\right].
 \label{erase:mixed-witness}
\end{equation}
Both numerator factors are positive on the stated intervals, so this expression has positive derivative in $r>0$. Two fixed admissible values of $\zeta$ therefore give different public directional gradients at leading order $\varrho\varepsilon^2$. The differentiated remainder is $O(\varrho\varepsilon^3+\varrho^2\varepsilon^2)$ on these compact sets. For sufficiently small positive gains, equal public data thus have different public derivatives while the gate is open. This expansion establishes nonidentity of the preclosure force; \cref{erase:separated-energy} is an exact postclosure identity.

For each selected pair $I$, retain its size and the empirical measure
\begin{equation}
 \eta_I=\frac1{|I|}\sum_{i\in I}\delta_{(x_i,y_i)},\qquad
 M=\sum_I|I|\eta_I.
 \label{erase:parent-state}
\end{equation}
This is a deterministic configuration statistic, not an assumed probability distribution over uncertain population configurations. Let $V[v;M]$ be the public Fisher vector field at the public state $v=(x,y)$, computed from $\mathcal E_{\rm pub}$. It is determined by $v$ and the counting measure $M$: source sums use $M-\delta_v$, and the derivative includes both receiver and incoming source terms from \cref{erase:public-derivatives}. The normalized Gibbs denominator is computed from the same counting measure. Permutation equivariance makes this definition independent of labels, including multiplicities.

\propositionheading{Exact parent evolution and public responses}{erase:parent-flow}
The inherited parent equations are
\begin{equation}
 \frac{d}{dt}\int f(v)\eta_I(dv)
 =\int\nabla f(v)\cdot V[v;M]\eta_I(dv)
 \label{erase:empirical-flow}
\end{equation}
for every continuously differentiable test function. On the realizable finite-atom states they determine the retained evolution uniquely. They also determine every declared fresh public response whose source selection depends only on public states, including joint responses with a shared sampled child.
\begin{proof}
Differentiate the finite sum in \cref{erase:parent-state} using the exact public equations. This gives \cref{erase:empirical-flow}. Equal retained parent measures admit enumerations that differ only by within-pair permutations. Uniqueness and permutation equivariance of the public system give identical evolved measures, proving the retained uniqueness assertion. For a fresh response, select the child or children using the declared normalized public kernel, then integrate their public sample laws against the retained empirical measures. Both the selection probabilities and sample laws are determined by the retained states. For shared-source outputs, use one shared draw rather than multiplying separate marginals; the same calculation preserves their joint law.
\end{proof}

The parent can carry belief and model marginals $q_I=\int b_x\eta_I(d(x,y))$ and $s_I=\int b_y\eta_I(d(x,y))$, with normalized public evaluator $G_J$. These two marginals alone are insufficient: $\eta_I$ retains the paired public roles and receiver-specific routing information. The inherited equations are \cref{erase:empirical-flow}, not a newly postulated Fisher flow for the two mixture marginals.

The physical compression here discards the private and gate laws: the admitted family has twenty-four continuous means, and the retained empirical state keeps eight public means, up to within-pair permutations. It keeps every public child state. Cached private readouts, feedback that changes the erased channels, and protocols that later inspect discarded private coordinates are outside the response class. The reference, gate, initial sign-box preparation, and boundary restrictions are all supplied assumptions. The dynamics select changing source preferences and a persistent grouping within them, but do not discover the erasing architecture. The two final sign boxes remain separated, so this example supplies no recursive second-stage formation and no further compression of the distinct public states.
